% Kartikey Pathak --- two-page detailed resume (for roles where depth helps,
% and as the version linked from the website)
% build:  pdflatex -interaction=nonstopmode resume-full.tex
\documentclass[a4paper,10pt]{article}
\usepackage[left=1.35cm,right=1.35cm,top=1.1cm,bottom=1.0cm]{geometry}

\newcommand{\secabove}{0.85em}
\newcommand{\secbelow}{0.45em}
\newcommand{\pointsep}{0.08em}
\newcommand{\pointafter}{0.25em}

% Shared preamble for both resume variants.
% Set \DENSE to 1 before % Shared preamble for both resume variants.
% Set \DENSE to 1 before % Shared preamble for both resume variants.
% Set \DENSE to 1 before \input{preamble} for the tighter one-page build.

\usepackage[T1]{fontenc}
\usepackage[utf8]{inputenc}
\usepackage{charter}
\usepackage{titlesec}
\usepackage{enumitem}
\usepackage{xcolor}
\usepackage[hidelinks]{hyperref}
\usepackage{tabularx}

\definecolor{ink}{HTML}{111111}
\definecolor{link}{HTML}{1A4B8C}
\color{ink}
\hypersetup{colorlinks=true,urlcolor=link,linkcolor=link}

\pagestyle{empty}
\setlength{\parindent}{0pt}
\raggedright

% ---- section headings -------------------------------------------------
\titleformat{\section}
  {\normalfont\large\bfseries\scshape}{}{0em}{}[\vspace{-0.72em}\rule{\linewidth}{0.6pt}]
\titlespacing{\section}{0pt}{\secabove}{\secbelow}

% ---- helpers ----------------------------------------------------------
% \entry{organisation}{dates}{role}{location}
\newcommand{\entry}[4]{%
  \begin{tabularx}{\linewidth}{@{}Xr@{}}
    \textbf{#1} & #2 \\
    \textit{#3} & \textit{#4} \\
  \end{tabularx}%
  \vspace{-0.15em}%
}
% \pentry{project}{stack}{dates}
\newcommand{\pentry}[3]{%
  \begin{tabularx}{\linewidth}{@{}Xr@{}}
    \textbf{#1} \,\textbar\, \textit{#2} & #3 \\
  \end{tabularx}%
  \vspace{-0.15em}%
}
\newenvironment{points}
  {\begin{itemize}[leftmargin=1.1em,itemsep=\pointsep,topsep=0.18em,parsep=0pt,label=--]}
  {\end{itemize}\vspace{\pointafter}}
\newcommand{\skillrow}[2]{\textbf{#1:} #2\\[0.12em]}

% ---- the header, identical in both variants ---------------------------
\newcommand{\resumeheader}{%
\begin{center}
  {\LARGE\bfseries Kartikey Pathak}\\[0.35em]
  {\small I build machines end to end --- mechanical, PCB, firmware, control, and the app on top.}\\[0.4em]
  {\small
    +91 7506048765 \,\textbar\,
    \href{mailto:kartikeypathak.imp@gmail.com}{kartikeypathak.imp@gmail.com} \,\textbar\,
    \href{https://spongefal.vercel.app}{spongefal.vercel.app} \,\textbar\,
    \href{https://github.com/NoobMaster-version}{github.com/NoobMaster-version} \,\textbar\,
    \href{https://www.linkedin.com/}{LinkedIn}
  }
\end{center}
}
 for the tighter one-page build.

\usepackage[T1]{fontenc}
\usepackage[utf8]{inputenc}
\usepackage{charter}
\usepackage{titlesec}
\usepackage{enumitem}
\usepackage{xcolor}
\usepackage[hidelinks]{hyperref}
\usepackage{tabularx}

\definecolor{ink}{HTML}{111111}
\definecolor{link}{HTML}{1A4B8C}
\color{ink}
\hypersetup{colorlinks=true,urlcolor=link,linkcolor=link}

\pagestyle{empty}
\setlength{\parindent}{0pt}
\raggedright

% ---- section headings -------------------------------------------------
\titleformat{\section}
  {\normalfont\large\bfseries\scshape}{}{0em}{}[\vspace{-0.72em}\rule{\linewidth}{0.6pt}]
\titlespacing{\section}{0pt}{\secabove}{\secbelow}

% ---- helpers ----------------------------------------------------------
% \entry{organisation}{dates}{role}{location}
\newcommand{\entry}[4]{%
  \begin{tabularx}{\linewidth}{@{}Xr@{}}
    \textbf{#1} & #2 \\
    \textit{#3} & \textit{#4} \\
  \end{tabularx}%
  \vspace{-0.15em}%
}
% \pentry{project}{stack}{dates}
\newcommand{\pentry}[3]{%
  \begin{tabularx}{\linewidth}{@{}Xr@{}}
    \textbf{#1} \,\textbar\, \textit{#2} & #3 \\
  \end{tabularx}%
  \vspace{-0.15em}%
}
\newenvironment{points}
  {\begin{itemize}[leftmargin=1.1em,itemsep=\pointsep,topsep=0.18em,parsep=0pt,label=--]}
  {\end{itemize}\vspace{\pointafter}}
\newcommand{\skillrow}[2]{\textbf{#1:} #2\\[0.12em]}

% ---- the header, identical in both variants ---------------------------
\newcommand{\resumeheader}{%
\begin{center}
  {\LARGE\bfseries Kartikey Pathak}\\[0.35em]
  {\small I build machines end to end --- mechanical, PCB, firmware, control, and the app on top.}\\[0.4em]
  {\small
    +91 7506048765 \,\textbar\,
    \href{mailto:kartikeypathak.imp@gmail.com}{kartikeypathak.imp@gmail.com} \,\textbar\,
    \href{https://spongefal.vercel.app}{spongefal.vercel.app} \,\textbar\,
    \href{https://github.com/NoobMaster-version}{github.com/NoobMaster-version} \,\textbar\,
    \href{https://www.linkedin.com/}{LinkedIn}
  }
\end{center}
}
 for the tighter one-page build.

\usepackage[T1]{fontenc}
\usepackage[utf8]{inputenc}
\usepackage{charter}
\usepackage{titlesec}
\usepackage{enumitem}
\usepackage{xcolor}
\usepackage[hidelinks]{hyperref}
\usepackage{tabularx}

\definecolor{ink}{HTML}{111111}
\definecolor{link}{HTML}{1A4B8C}
\color{ink}
\hypersetup{colorlinks=true,urlcolor=link,linkcolor=link}

\pagestyle{empty}
\setlength{\parindent}{0pt}
\raggedright

% ---- section headings -------------------------------------------------
\titleformat{\section}
  {\normalfont\large\bfseries\scshape}{}{0em}{}[\vspace{-0.72em}\rule{\linewidth}{0.6pt}]
\titlespacing{\section}{0pt}{\secabove}{\secbelow}

% ---- helpers ----------------------------------------------------------
% \entry{organisation}{dates}{role}{location}
\newcommand{\entry}[4]{%
  \begin{tabularx}{\linewidth}{@{}Xr@{}}
    \textbf{#1} & #2 \\
    \textit{#3} & \textit{#4} \\
  \end{tabularx}%
  \vspace{-0.15em}%
}
% \pentry{project}{stack}{dates}
\newcommand{\pentry}[3]{%
  \begin{tabularx}{\linewidth}{@{}Xr@{}}
    \textbf{#1} \,\textbar\, \textit{#2} & #3 \\
  \end{tabularx}%
  \vspace{-0.15em}%
}
\newenvironment{points}
  {\begin{itemize}[leftmargin=1.1em,itemsep=\pointsep,topsep=0.18em,parsep=0pt,label=--]}
  {\end{itemize}\vspace{\pointafter}}
\newcommand{\skillrow}[2]{\textbf{#1:} #2\\[0.12em]}

% ---- the header, identical in both variants ---------------------------
\newcommand{\resumeheader}{%
\begin{center}
  {\LARGE\bfseries Kartikey Pathak}\\[0.35em]
  {\small I build machines end to end --- mechanical, PCB, firmware, control, and the app on top.}\\[0.4em]
  {\small
    +91 7506048765 \,\textbar\,
    \href{mailto:kartikeypathak.imp@gmail.com}{kartikeypathak.imp@gmail.com} \,\textbar\,
    \href{https://spongefal.vercel.app}{spongefal.vercel.app} \,\textbar\,
    \href{https://github.com/NoobMaster-version}{github.com/NoobMaster-version} \,\textbar\,
    \href{https://www.linkedin.com/}{LinkedIn}
  }
\end{center}
}


\begin{document}
\resumeheader

% =======================================================================
\section{Education}
\entry{Veermata Jijabai Technological Institute (VJTI)}{Mumbai, India}
      {B.Tech, Electrical Engineering --- Minor in Robotics (CGPA 8.04)}{2023 -- 2027}
\vspace{0.35em}
\entry{Matoshree Prabodhini Jr.\ College}{Maharashtra, India}
      {HSC, Maharashtra State Board --- MHT-CET 99.14 percentile}{2023}

% =======================================================================
\section{Experience}

\entry{Eyecandy Robotics}{Nov 2025 -- Present}
      {Electronics / Hardware Engineer}{}
\begin{points}
  \item Design the PCB stack for a legged companion robot: compute boards (Raspberry Pi Zero 2W and Milk-V SoC carriers), 1S/2S power and charging boards (BQ25887, MP2672, TPS562202 buck rails), a servo motor-driver board, a microphone board and inter-board connector boards --- schematic through layout, DFM and fabrication output.
  \item Board-level integration of an MPU6050 IMU, a Feetech servo bus with logic buffering, per-leg force-sensitive-resistor conditioning (LM393), PAM8403 audio and a ReSpeaker mic array; laid out test points and solder-jumper options so each subsystem can be probed or bypassed during bring-up.
  \item Wrote the fleet-provisioning tooling that clones one built Raspberry Pi runtime image onto many robots with unique machine IDs, hostnames and SSH host keys, so a fleet coexists on one network without IP or config collisions.
\end{points}

\entry{NishCorp Technologies}{May -- Nov 2025}
      {Robotics Control Systems Intern}{Hybrid}
\begin{points}
  \item Turned the CAD of a stair-climbing wheelchair into a working simulation: CAD to URDF, spawned in Gazebo/RViz, differential-drive kinematics, IMU plugin and teleop --- so motion and steering could be studied before anything was built in steel.
  \item Helped build the half-scale hardware prototype through 3D printing, assembly, motor electronics and field testing; it could sit and stand under its own power.
  \item The simulation work formed part of the technical case that helped the startup raise Rs 13 lakh in grants.
\end{points}

\entry{Krishna Defence and Allied Industries Ltd.}{May -- Jul 2025}
      {Research Intern, Defence Technology}{at IIT Gandhinagar}
\begin{points}
  \item Solved indoor localization for an autonomous forklift in a naval dockyard, where GPS is blocked by steel and wheel odometry drifts; surveyed Bluetooth, LoRa, RFID, radar and vision before selecting ultra-wideband ranging.
  \item Wrote custom double-sided two-way ranging firmware for the Qorvo DWM3001CDK, tuning channel, preamble length and PAC size and re-deriving timing delays per configuration --- extending usable range from the stock 45 m to \textbf{315 m}.
  \item Ran SLAM and NAV2 on a Raspberry Pi 4B fusing UWB, IMU and LIDAR; scaled testing from a $2\times2$ m room to a $50\times60$ m yard and validated against a surveyor's total station rather than our own code: \textbf{0.4\% error over 10,000 m\textsuperscript{2}}.
\end{points}

% =======================================================================
\section{Projects}

\pentry{PCB Axial-Flux BLDC Motor}{KiCAD, FEMM, Python, STM32, FOC --- final year project}{2026 -- Present}
\begin{points}
  \item Designing a coreless axial-flux motor whose stator \emph{is} the PCB: 100 mm 2-layer board, 6 wedge coils over 8 poles, dual 3D-printed rotor carrying 16 N42 magnets across a 0.5 mm air gap. Coil geometry is generated programmatically rather than drawn by hand.
  \item Built the magnetic model myself after establishing that a top-down 2D cut cannot yield torque: unrolled the machine on a cylinder into five radial slices to put coil current out of plane, then extracted Kt, Kv, back-EMF and ripple from the solved air-gap flux. Predicted 40.9 mN$\cdot$m/A rms, 496 rpm/V and zero cogging torque.
  \item Simulation caught a design error before fabrication --- the plastic rotor gives away $3.3\times$ the torque of the intended steel-backed design --- and showed the motor is thermally, not magnetically, limited. Sized 6 mm steel back plates as the fix by checking plate saturation.
  \item Companion driver firmware on an STM32G431 with a DRV8311H: centre-aligned 3-phase PWM at 20 kHz, a 40 kHz control ISR with zero-sequence injection, AS5600 magnetic encoder over I2C and DMA-fed phase current sensing.
\end{points}

\pentry{Smart Switch Board}{ESP32, ESP-IDF, FreeRTOS, MQTT/TLS, BLE, Flutter, CAD}{2025 -- Present}
\begin{points}
  \item Built an internet-controlled mains wall switch across every layer alone: CAD enclosure, a relay and optocoupler board isolating 220 V from a 3.3 V MCU, the firmware, a TLS MQTT broker and a Flutter app.
  \item Held to a hard rule that the physical switches keep working with WiFi, cloud and app all down --- WiFi provisioning over BLE with nothing hardcoded, state persisted in NVS across power cuts, and a 5-second heartbeat keeping switch, firmware, broker and app agreed on state. Installed and running at home.
\end{points}

\pentry{e-Yantra Warehouse Drone, IIT Bombay}{ROS2, OpenCV, ArUco, PID, path planning}{Oct 2024 -- Jan 2025}
\begin{points}
  \item Tuned PID control for stable hover and precise waypoint motion, and wrote the vision that detects and decodes ArUco markers for localization and pick-and-place in a simulated warehouse. \textbf{National semifinalist.}
\end{points}

\pentry{MARIO --- 3-DOF Teaching Arm}{ROS2, Gazebo Fortress, micro-ROS, ESP32, URDF, C++}{Jul -- Aug 2024}
\begin{points}
  \item Ported my college robotics society's teaching arm off end-of-life Gazebo Classic --- model, plugins and launch files --- rewrote teleop for beginners, and unified the ESP32/micro-ROS pipeline so one command set drives both simulation and hardware. Still runs the ROS2 workshops.
\end{points}

\pentry{Also built}{vision/ML, embedded, mechanical}{2023 -- 2025}
\begin{points}
  \item \textbf{Coin sorter} --- 24-hour hackathon; a camera plus a transfer-learned CNN on a Raspberry Pi classifies Rs 1/2/5/10 and servo gates route each coin, above 95\% on Rs 5 and Rs 10.
  \item \textbf{Voice-controlled pick-and-place arm} --- speech recognition to object detection to inverse kinematics in ROS2/Gazebo. \textbf{Kurma} --- statically stable quadruped with a turtle-derived gait, Mass Robotics entrant.
  \item \textbf{Self-balancing robot} --- complementary-filter sensor fusion and PID on ESP-IDF; the society now teaches on kits of it. \textbf{Maze solver} --- IR array, 100 Hz loop, left-hand-rule mapping with shortest-path replay, 4th place. \textbf{Gesture-controlled car} --- glove IMU over ESP-NOW. Plus a cat feeder and a river-cleaning boat.
\end{points}

% =======================================================================
\section{Teaching \& Leadership}

\entry{Society of Robotics and Automation, VJTI}{Aug 2024 -- Mar 2025}
      {Active Member --- workshops and mentoring}{VJTI, Mumbai}
\begin{points}
  \item Taught workshops on embedded C, ESP32 and sensor fusion (MPU6050, complementary filters) and mentored junior teams through their first self-balancing and line-following robots. The society manufactures and assembles the robot kits these workshops run on; \textbf{over 200 students} have learned on them.
\end{points}

\entry{Pratibimb, VJTI Cultural Festival}{Aug 2023 -- Mar 2025}
      {Department Head, Electrical Engineering}{VJTI, Mumbai}
\begin{points}
  \item Led the Electrical Engineering contingent to its first overall win in the competition's 26-year history. Separately directed a short film, \textit{Plot Se Panga}, which took first prize across all branches.
\end{points}

% =======================================================================
\section{Technical Skills}
\skillrow{Languages}{C, C++, Python, Dart, Lua, TypeScript}
\skillrow{Robotics}{ROS \& ROS2, Gazebo (Classic \& Fortress), RViz, URDF, micro-ROS, SLAM, NAV2}
\skillrow{Control}{PID tuning, sensor fusion, forward \& inverse kinematics, gait design, field-oriented control}
\skillrow{Embedded}{ESP32 (ESP-IDF, FreeRTOS, ESP-NOW), STM32, Raspberry Pi, ARM Cortex-M, I2C/SPI/UART, PWM, ADC/DMA, BLE, MQTT}
\skillrow{Electronics}{KiCAD schematic \& layout, DFM and fabrication output, buck regulators, battery charging, motor drivers, mains isolation, soldering and bring-up with scope and logic analyzer}
\skillrow{Simulation \& analysis}{FEMM magnetostatics, Gazebo physics, OnShape CAD, 3D printing}
\skillrow{Vision \& ML}{OpenCV, ArUco, TensorFlow, CNNs and transfer learning}
\skillrow{Tools}{Linux, Git, shell scripting, Astro (built my own portfolio site)}

\end{document}
